% ---
% title: "Constructing Affine Invariant/Equivariant MCMC Samplers"
% date: 2026-03-18
% ---
\documentclass{article}
\usepackage{amsmath, amsfonts, amssymb} % Math formatting
\usepackage{graphicx}

\usepackage{mdframed} % Framed Enviroments
\newmdenv[ 
  topline=false,
  bottomline=false,
  skipabove=\topsep,
  skipbelow=\topsep
]{sidework} %% Side-work

\title{Constructing Affine Invariant/Equivariant MCMC Samplers}
\date{}

\begin{document}

\maketitle

Affine invariance/equivariance is a desirable property for statistical inference algorithms. For example, consider sampling the posterior $p(\theta | \mathcal D)$, where $\theta = (\text{Lenght}, \text{Mass}, ...)$. If the data $\mathcal D$ forces the posterior to different length/mass scales, this will correspond to stretchin/squishing the probability distribution (which are affine transformations), hence an affine invariant algorithm could prevent some heart-ache here.

For MCMC, let's be specific, what is actually conserved under affine transformations? For a fixed RNG seed, the Metroplist-Hastings ratio is invariant under pushforward of the target distribution. This means the walkers' positions are equivariant with respect to pushforwards of the target distribution.



\section{Stretch, Walk, \& Side Move}

\section{Gradient Based Samplers}


\subsection{Langevin}

\subsection{Hamiltonian Monte Carlo}






\end{document}
